\section{Datos públicos}\label{sec:externos}
Las bases permiten agregar información pública. Con ella se podrían buscar dos cosas, más rendimientos para entrenar
o más información sobre las mismas parcelas.

\textbf{Más rendimientos.} No encontramos ningún conjunto público con rendimiento de cebada por parcela en la zona
(Tabla~\ref{tab:externos}). Los 138 rendimientos entregados son todo lo que hay a esa escala. Dentro del propio
conjunto, Sentinel-2 y Landsat cubren de 2022 a 2024 las mismas 197 parcelas. No traen rendimiento, pero muestran
cómo se comporta cada parcela año con año, y eso se puede usar como variable.

\begin{table}[!t]
\caption{Fuentes revisadas como posibles rendimientos}\label{tab:externos}
\centering\footnotesize
\begin{tabularx}{\columnwidth}{@{}>{\raggedright\arraybackslash}p{2.4cm}X@{}}
\toprule
Fuente & Motivo por el que no sirve \\
\midrule
SIAP, producción municipal~\cite{siap2026} & Es un promedio por municipio y va en sentido contrario a las etiquetas \\
Ensayos de variedades del CIMMYT en Valles Altos & Son parcelas experimentales de otros años y en parte con riego \\
Bitácoras de productores MasAgro & No se publican a nivel de parcela \\
\bottomrule
\end{tabularx}
\end{table}

\textbf{Más información sobre las parcelas.} Ordenamos las fuentes según lo que esperamos que aporten.
\begin{enumerate}
  \item Radar de Sentinel-1, que atraviesa las nubes y puede detectar acame y fecha de cosecha~\cite{chauhan2020}.
  \item Lluvia y temperatura diarias, para medir lo que pasó cerca de la floración de cada parcela.
  \item Bandas originales de Sentinel-2, con máscara de nubes por píxel y sin los bordes de cada parcela, que en las
        parcelas pequeñas pesan mucho.
  \item Propiedades del suelo y un modelo de elevación más fino.
  \item Representaciones anuales de imágenes satelitales publicadas por Google, de 2017 a 2024, como resumen del
        historial de cada sitio.
\end{enumerate}
